% ===== §2 =====
\section{The Being of a Relation}
\label{sec:existence}

\noindent
This section takes the first of the six dimensions, the sheer existence of a relation, and establishes a single result: that the being of a bond, its \emph{that it is} as distinct from any account of what it is like, cannot be grounded by any proposition, and that this ungroundability is not a defect in relational language but the condition under which the existence of a relation can be a matter of fidelity rather than of record. The section states the phenomenon, sets it against the two traditions that would seem to close it, isolates the core that neither reaches, and draws the claim.

\subsection*{The phenomenon}

Take first the plain fact that a relation exists. Two people are bound, or they are not; the bond is alive, or it has died. Language seems obliged to be able to say which, and it holds the sentences ready to hand. We are together. This is real. I am yours. These sentences present themselves as if they could establish the existence they announce, as if to say the bond were in some measure to make it, or at least to fix and secure it against doubt. The appearance is strong enough that a whole practice of reassurance rests on it, the practice of asking to be told, and of telling, that the bond is there.

They do not establish it. The sentence \emph{we are together} neither proves that the bond exists nor brings it into being. It can be spoken over a bond already hollow, and it can be withheld from a bond fully alive. The existence of a relation is not the kind of thing a proposition can ground. Whether the bond is real is settled somewhere the sentence does not reach, and every sentence about it arrives after the fact it would establish, a report on a state it did not produce and cannot guarantee. The being of the relation lies outside all the signifiers that would name it.

\subsection*{Two ways of trying to close the gap, and why they do not}

Two traditions promise to close exactly this gap, and it is worth seeing why neither does, since the reason is instructive.

The first is the performative account of speech. Austin taught that some utterances do not describe a state of affairs but bring one about, that to say the words is, under the right conditions, to do the deed \citep{austin1962words}. I now pronounce you; I promise; I name this ship. It can look as though \emph{I am yours} belongs to this class, an utterance that makes the bond it declares. But the performative works only where a convention stands ready to be invoked and a procedure exists whose completion the words execute. The marrying words marry because an institution recognizes them; the naming names because a practice confers the name. The existence of a living bond has no such institution behind the words. There is no procedure whose correct performance constitutes the reality of love, no convention on whose authority \emph{I am yours} could execute the bond it speaks. The performative closes the gap between word and deed only where an institution has already closed it; for the bare existence of a relation, no institution has, and the words fall back into report.

The second is the ideal of transparent communication, the thought that a bond is secured when the two have made themselves fully intelligible to each other, when nothing of either remains unsaid. Habermas built a whole ethics on the regulative ideal of communication free from distortion, oriented to mutual understanding \citep{habermas1984theory}. On such a view the existence of a good bond would be underwritten by the completeness of the understanding reached within it. But understanding, however complete, is understanding of what the other says and shows, and the existence of the bond is not something either party says or shows. It is presupposed by everything they say and shown by nothing. Two people can understand each other to the last word and still face the question of whether the bond between them is alive, because that question is not answered at the level of understanding at all. The transparency ideal, pressed to its limit, would deliver a perfect mutual intelligibility beneath which the existence of the bond would remain exactly as ungrounded as before.

\subsection*{The core}

What both traditions reach for and miss is the same. The core language cannot reach here is the sheer existence of the bond, its \emph{that it is}, as against all the descriptions of what it is like. Descriptions there can be in abundance. The tone of a morning, the history of a quarrel and its repair, the shape of an ordinary evening, all of this can be said, and much of it can be said well. What cannot be said, because it is of a different order from any description, is the fact that beneath all these there is a bond at all, living now, in force now. That fact is presupposed by every description and secured by none of them. It is not a further, especially deep description that better language might supply. It is not a description at all, and that is why no description reaches it.

\begin{claim}[The ungroundability of the bond]
The existence of a relation, its being-in-force as distinct from every account of its qualities, cannot be established by any proposition. No performative constitutes it, for it invokes no institution; no completeness of mutual understanding underwrites it, for it is presupposed by all understanding and shown by none. The being of the bond lies outside the signifier that would name it.
\end{claim}

\begin{casebox}{the reassurance that settles nothing}
One partner asks to be told that the bond is real, and the other tells them, and the words are true and warmly meant. Yet the asking returns, not from distrust of the answer but because the answer, however true, is not the kind of thing that could settle the question. What was wanted was not a sentence but the standing of the bond itself, and no sentence delivers that standing. The recurrence of the asking is not a fault in either party. It is the mark that the existence of a bond lies where no reassurance can put it, outside the reach of any words that would confirm it.
\end{casebox}

\subsection*{Why the limit is generative here}

Far from being a defect, this ungroundability is what lets the existence of a relation be a matter of fidelity rather than of record. Because no sentence can ground the bond, the bond rests on something other than a sentence, on a holding that must be renewed and cannot be filed away as done. Consider the counter-case, a relation whose existence \emph{could} be established by a form of words. Its existence, once the words were produced, would be settled, on record, no longer at stake; and a bond no longer at stake is a bond one possesses by having produced its title. The limit forbids exactly this. Because the existence of the relation can never be put on record, it can never be possessed as a title, and so it remains what it must be to be a bond at all, a thing continually at stake, held open, renewed in the holding. The reader will recognize, without its needing to be argued here, that a promise carrying no external guarantee stands closer to this condition than a contract does: the contract puts the bond on record and thereby makes it a title; the promise, guaranteeing nothing, leaves the bond where the existence of a relation always is, outside the record and in the keeping. The ungroundability that looked like a poverty of language is the very thing that keeps the existence of a relation an act.
